\section*{Capítulo \ref{ch:g12:radioactive-decay} --- Decaimento radioativo}

\begin{solution}{exo:g12:radioactive-decay:1}
$T_{1/2} = \ln 2/\lambda = 0.693/\num{1.0e-6} = \qty{6.9e5}{s} \approx
8.0$ dias --- o relógio do iodo-131.
\end{solution}

\begin{solution}{exo:g12:radioactive-decay:2}
$T_{1/2} = 3.8 \times \num{86400} = \qty{3.28e5}{s}$;
$\lambda = 0.693/\num{3.28e5} = \qty{2.1e-6}{s^{-1}}$;
$\mathcal A = \lambda N = \num{2.1e-6} \times \num{1.0e10} \approx
\qty{2.1e4}{Bq}$.
\end{solution}

\begin{solution}{exo:g12:radioactive-decay:3}
$\dd N/\dd t = N_0(-\lambda)e^{-\lambda t} = -\lambda N$ --- a equação
vale; $N(0) = N_0 e^0 = N_0$;
$N(T_{1/2})/N_0 = e^{-\lambda \ln 2/\lambda} = e^{-\ln 2} = \tfrac12$.
\end{solution}

\begin{solution}{exo:g12:radioactive-decay:4}
$90$ anos $= 3\,T_{1/2}$: $\mathcal A = \num{4.0e5}/2^3 =
\qty{5.0e4}{Bq}$. $300$ anos $= 10\,T_{1/2}$:
$\mathcal A = \num{4.0e5}/1024 \approx \qty{3.9e2}{Bq}$.
\end{solution}

\begin{solution}{exo:g12:radioactive-decay:5}
Inclinação $= (18.85 - 20.00)/10 = -0.115$, logo
$\lambda = \qty{0.115}{h^{-1}} = 0.115/3600 = \qty{3.2e-5}{s^{-1}}$;
$T_{1/2} = 0.693/0.115 = 6.0$ horas --- o tecnécio-99m.
\end{solution}

\begin{solution}{exo:g12:radioactive-decay:6}
$N = (1.0/226) \times \num{6.02e23} = \num{2.7e21}$;
$T_{1/2} = 1600 \times \num{3.156e7} = \qty{5.05e10}{s}$, logo
$\lambda = \qty{1.37e-11}{s^{-1}}$;
$\mathcal A = \lambda N \approx \qty{3.7e10}{Bq}$ --- um curie, a
\omterm{def:g11:nucleus-radioactivity:activity}{atividade} do grama de rádio de Marie Curie.
\end{solution}

\begin{solution}{exo:g12:radioactive-decay:7}
Os dois esperam que sobre metade: $10$, e \num{1.0e20}. As flutuações são
$\sim 1/\sqrt{n}$ da contagem: com $10$ decaimentos, cerca de $30\%$ ---
$7$ ou $13$ não seriam surpresa; com \num{1.0e20}, cerca de $10^{-10}$.
Só a amostra grande obedece à lei com precisão mensurável.
\end{solution}

\begin{solution}{exo:g12:radioactive-decay:8}
$m = 140 \times \num{1.2e-4} = \qty{1.7e-2}{g}$;
$N = (0.0168/40) \times \num{6.02e23} = \num{2.5e20}$.
$\lambda = 0.693/(\num{1.25e9} \times \num{3.156e7}) =
\qty{1.76e-17}{s^{-1}}$;
$\mathcal A = \lambda N \approx \qty{4.4e3}{Bq}$ --- cerca de metade do
total do corpo (o grosso do resto é carbono-14).
\end{solution}

\begin{solution}{exo:g12:radioactive-decay:9}
$\lambda = 0.693/6.0 = \qty{0.116}{h^{-1}}$. Em $15$ horas:
$\mathcal A = 500\,e^{-0.116 \times 15} = 500 \times 0.177 \approx
\qty{88}{MBq}$. Abaixo de \qty{1}{MBq}:
$t = \ln(500)/0.116 \approx \qty{54}{horas}$ --- cerca de $9$
meias-vidas.
\end{solution}

\begin{solution}{exo:g12:radioactive-decay:10}
Carvão: $3.4/13.6 = \tfrac14$, exatamente $2$ meias-vidas:
$t = \num{1.15e4}$ anos. Osso:
$t = \num{5730} \times \ln(1/0.30)/\ln 2 = \num{5730} \times 1.74
\approx \num{1.0e4}$ anos.
\end{solution}

\begin{solution}{exo:g12:radioactive-decay:11}
$\lambda = 0.693/3.8 = \qty{0.182}{d^{-1}}$;
$t = \ln(100)/\lambda = 4.61/0.182 \approx 25$ dias (cerca de $6.6$
meias-vidas). Mas a cadeia do urânio no granito e no solo em volta gera
radônio continuamente: um arejamento único é desfeito em semanas, e por
isso a ventilação precisa ser permanente.
\end{solution}

\begin{solution}{exo:g12:radioactive-decay:12}
$N/N_0 = e^{-\lambda t}$ dá $\lambda t = \ln(N_0/N)$, e
$\lambda = \ln 2/T_{1/2}$ dá $t = T_{1/2}\ln(N_0/N)/\ln 2$. Aqui
$t = \num{5730} \times \ln(20)/\ln 2 = \num{5730} \times 4.32 \approx
\num{2.5e4}$ anos.
\end{solution}

\begin{solution}{exo:g12:radioactive-decay:13}
$N_{\mathrm K} = N_0 e^{-\lambda t}$ e
$N_{\mathrm{Ar}} = N_0 - N_{\mathrm K}$, logo
$N_{\mathrm{Ar}}/N_{\mathrm K} = e^{\lambda t} - 1$. Razão $3.0$:
$e^{\lambda t} = 4$, logo $\lambda t = 2\ln 2$ e
$t = 2\,T_{1/2} = \num{2.5e9}$ anos.
\end{solution}

\begin{solution}{exo:g12:radioactive-decay:14}
$\ln \mathcal A = 5.99$, $5.42$, $4.84$, $4.26$, $3.69$: quedas de
$0.57$--$0.58$ a cada \qty{60}{s} --- alinhados. Pela inclinação:
$\lambda = 2.30/240 = \qty{9.6e-3}{s^{-1}}$;
$T_{1/2} = 0.693/\num{9.6e-3} \approx \qty{72}{s}$.
\end{solution}

\begin{solution}{exo:g12:radioactive-decay:15}
(a) $2^{-10} \approx \num{1e-3}$. (b) $(1/2)^8 = 1/256 \approx 0.4\%$.
(c) Esperam-se $1000/8 = 125$ --- mas a lei se cala sobre quais: ela
prevê populações com exatidão e indivíduos de modo algum.
\end{solution}

\begin{solution}{pb:g12:radioactive-decay:1}
\textbf{1.} A troca com a atmosfera (comer, respirar) mantém a proporção
completada no valor atmosférico; a morte interrompe a entrada e o
decaimento assume.

\textbf{2.} $T_{1/2} = \num{5730} \times \num{3.156e7} =
\qty{1.81e11}{s}$, logo $\lambda = 0.693/\num{1.81e11} =
\qty{3.8e-12}{s^{-1}}$.

\textbf{3.} $\dd N/\dd t = -\lambda N_0 e^{-\lambda t} = -\lambda N$, e
$N(0) = N_0$.

\textbf{4.} $\mathcal A_0 = 13.6/60 = \qty{0.227}{Bq}$ por grama;
$N = \mathcal A_0/\lambda = 0.227/\num{3.8e-12} \approx \num{5.9e10}$
\omterm{def:g11:fundamental-interactions:ladder}{átomos} por grama.

\textbf{5.} $(1/12) \times \num{6.02e23} = \num{5.0e22}$ \omterm{def:g11:fundamental-interactions:ladder}{átomos} de
carbono; proporção $\num{5.9e10}/\num{5.0e22} \approx \num{1.2e-12}$ ---
um em $10^{12}$, como prometido.

\textbf{6.} $\mathcal A = \mathcal A_0 e^{-\lambda t}$ dá
$\lambda t = \ln(\mathcal A_0/\mathcal A)$, e
$\lambda = \ln 2/T_{1/2}$ transforma isso em
$t = T_{1/2}\ln(\mathcal A_0/\mathcal A)/\ln 2$.

\textbf{7.} $t = \num{5730} \times \ln(13.6/9.1)/\ln 2 =
\num{5730} \times 0.580 \approx \num{3.3e3}$ anos.

\textbf{8.} A morte da madeira --- o corte da árvore, e não o entalhe;
uma estátua feita de madeira velha (ou uma falsificação entalhada em
madeira antiga) é anterior ou posterior à data do material dela.

\textbf{9.} $t = \num{5730} \times \ln(13.6/10.4)/\ln 2 =
\num{5730} \times 0.387 \approx \num{2.2e3}$ anos --- plenamente
compatível com um escriba de vinte e dois séculos atrás.

\textbf{10.} $\Delta t \approx \Delta\mathcal A/(\lambda\mathcal A)$ com
$\lambda = \ln 2/5730 = \num{1.21e-4}$ por ano:
$\Delta t = (0.2/10.4)/\num{1.21e-4} \approx \num{1.6e2}$ anos --- uma
data boa até um século ou dois.

\textbf{11.} $t = \num{5730} \times \ln(13.6/13.4)/\ln 2 \approx 120$
anos --- indistinguível de moderno dentro da incerteza de $\pm 160$
anos: o pergaminho é recente, e a ``antiguidade'', falsa.

\textbf{12.} $13.6/2^n < 0.2$ exige $2^n > 68$: $n = 7$ meias-vidas
(pois $2^7 = 128$), cerca de \num{4.0e4} anos.

\textbf{13.} Dez meias-vidas: cerca de \num{5.7e4} anos --- o horizonte
prático, uns sessenta mil anos.

\textbf{14.} $\num{6.6e7}/\num{5730} \approx \num{11500}$ meias-vidas.
Os \num{5.9e10} \omterm{def:g11:fundamental-interactions:ladder}{átomos} de um grama se esgotam quando
$2^n > \num{5.9e10}$, ou seja, $n \approx 36$ meias-vidas
($\approx \num{2e5}$ anos): muito antes da idade dos dinossauros, não
resta um único \omterm{def:g11:fundamental-interactions:ladder}{átomo} de carbono-14 --- não há nada para contar.

\textbf{15.} Um relógio bem mais lento: potássio-40 ou urânio-238, com
meias-vidas de bilhões de anos.

\textbf{16.} A lava derretida deixa o argônio borbulhar embora, de modo
que o cristal recém-solidificado guarda potássio e zero argônio: a razão
começa em $0$ na solidificação --- o relógio está zerado.

\textbf{17.} Como no exercício: $N_{\mathrm K} = N_0 e^{-\lambda t}$ e
$N_{\mathrm{Ar}} = N_0 - N_{\mathrm K}$, logo
$N_{\mathrm{Ar}}/N_{\mathrm K} = e^{\lambda t} - 1$.

\textbf{18.} $e^{\lambda t} = 1.60$:
$t = \num{1.25e9} \times \ln(1.60)/\ln 2 = \num{1.25e9} \times 0.678
\approx \num{8.5e8}$ anos.

\textbf{19.} Não, nos dois sentidos: em \num{8.5e8} anos o carbono-14 já
correu $10^5$ meias-vidas --- silêncio; em \num{3.3e3} anos o
potássio-40 correu \num{2.6e-6} de uma \omterm{def:g11:nucleus-radioactivity:halflife}{meia-vida} --- a razão de argônio
$\approx \lambda t \sim \num{2e-6}$ é pequena demais para ser medida.
Cada relógio só lê de cerca de um décimo de \omterm{def:g11:nucleus-radioactivity:halflife}{meia-vida} até cerca de dez.

\textbf{20.} Estátua $\approx \num{3.3e3}$ anos, papiro
$\approx \num{2.2e3}$ anos, rocha $\approx \num{8.5e8}$ anos --- três
relógios, uma lei: $N = N_0\,e^{-\lambda t}$, lida ao contrário.
\end{solution}
